% ======================================================================
% SECTION 4: AUTONOMOUS CLINICAL OPERATIONS AND SITE MANAGEMENT
% File: sections/clinical_operations.tex
% ======================================================================

[PROCESSING INSTRUCTION: Generate this section (approximately 3-4 pages) covering autonomous
clinical operations that replace the traditional ClinOps team. Use the following source files:

SOURCE FILES TO PROCESS:
- sponsor/input\_files/sponsor\_04\_operations\_enrollment.md (SSU delays, site selection, enrollment,
  decentralized elements, diversity requirements, failure mode mitigation)
- sponsor/input\_files/org\_04\_functions\_site\_enrollment.md (RWD-driven feasibility, CRO-led
  execution, hybrid model, centralized analytics, LPLV-to-lock cycle times)
- sponsor/input\_files/org\_02\_operating\_model\_roles.md (Clinical operations staffing model)
- national-platform/new\_paper/final\_paper/sections/site\_establishment.tex (site documentation)
- national-platform/new\_paper/final\_paper/sections/patient\_journey.tex (patient pipeline)

SUBSECTION STRUCTURE:

\subsection{Study Orchestrator Agent (study\_orchestrator)}
- End-to-end study plan management: dependencies, work queues, issue routing
- Maps to: Cross-functional study team coordination
- Integration with all other agents through event-driven messaging

\subsection{ClinOps Agent (clinops\_agent)}
- Site startup automation: site qualification, training, regulatory package submission
- Monitoring plan generation: risk-based monitoring per ICH E6(R3)
- Deviation tracking and CAPA generation
- Enrollment tracking with predictive analytics
- Query resolution workflows

\subsection{Enrollment Optimization and Diversity}
- Central lab coordination, EHR matching, prescreening networks
- Decentralized trial elements: telehealth, local labs, home nursing, DHTs
- FDORA diversity action plan automation
- Screen failure prediction and mitigation

\subsection{Failure Mode Mitigation}
- 5-row failure mode table from sponsor\_04: SSU delays, high screen failure,
  protocol amendments, dropouts, supply interruptions
- Autonomous detection and response algorithms
- Gantt timeline management from TPP through launch

PYTHON SCRIPT REQUIREMENTS (to be generated by Claude Code Opus 4.6 in the future):
- File: sponsor/template/scripts/study\_orchestrator.py
  Purpose: Master orchestrator managing study plan, dependencies, and work queues.
  Implements event-driven architecture routing tasks to specialized agents.
  Input: study config, agent registry. Output: task assignments, status updates.
- File: sponsor/template/scripts/clinops\_agent.py
  Purpose: Site startup automation, monitoring plan generation, enrollment tracking.
  Implements RBM per ICH E6(R3) with centralized statistical monitoring.
  Include SSU cycle time tracking against median 61-120 day benchmarks.
- File: sponsor/template/scripts/enrollment\_optimizer.py
  Purpose: Enrollment prediction engine with diversity tracking. Parse EHR data
  for pre-screening, compute enrollment velocity, generate diversity reports
  per FDORA requirements. Include screen failure prediction model.
- File: sponsor/template/scripts/failure\_mode\_detector.py
  Purpose: Real-time detection of the 5 failure modes (SSU delays, screen failure,
  protocol amendments, dropouts, supply interruptions). Trigger automated
  mitigation responses and escalation alerts.
- All scripts must pass ruff lint (line-length 120, E/F/W rules).
- Process with Claude Code Opus 4.6 (1M token context).

TABLE REQUIREMENTS:
- Table 5: Study Startup Automation Workflow
  Columns: Step | Traditional Timeline | Autonomous Timeline | Agent Owner | Key Automation
- Table 6: Failure Mode Mitigation Matrix (5 modes from sponsor\_04)
  Columns: Failure Mode | Detection Signal | Autonomous Response | Escalation Trigger | KPI

FORMATTING NOTES:
- Include the Gantt-style timeline as a text-based diagram.
- Reference adapted ICH E6(R3) centralized monitoring provisions.]

\subsection{Study Orchestrator Agent}

[PROCESSING INSTRUCTION: Generate 1 page on the study\_orchestrator. This is the central
coordination agent that replaces the cross-functional study team. It manages the end-to-end
study plan, maintains dependency graphs between agents, routes work items, and escalates
issues. Describe the event-driven architecture: each agent publishes status events, and
the orchestrator consumes them to maintain global study state. Reference the study/trial
layer from org\_02\_operating\_model\_roles.md. Describe integration with the 7 governance
forums (Portfolio Committee, Study Steering Committee, DMC/IDMC, Safety Management Team,
Dose Escalation Committee, Quality/Risk Review Board, Vendor Governance).
Process with Claude Code Opus 4.6 (1M token context).]

\subsection{Clinical Operations Agent}

[PROCESSING INSTRUCTION: Generate 1.5 pages on the clinops\_agent. Cover: site startup
automation (site qualification scoring, training module delivery, regulatory package
assembly and submission), monitoring plan generation using risk-based monitoring per
ICH E6(R3) with QTLs/KRIs, deviation tracking with automated CAPA generation, and
query resolution. Reference the operational failure root causes from sponsor\_04:
SSU delays (61-120+ day timelines), NCI vs. industry gap. Reference the hybrid
sponsor-CRO model from org\_04. Include site selection using RWD-driven feasibility.
Process with Claude Code Opus 4.6 (1M token context).]

\subsection{Enrollment Optimization and Diversity}

[PROCESSING INSTRUCTION: Generate 1 page on enrollment optimization. Cover: central lab
coordination, EHR matching for pre-screening, decentralized trial elements (telehealth,
local labs, home nursing, digital health technologies), and FDORA diversity action plan
automation. Reference the enrollment ecosystem from sponsor\_04 and the diversity
requirements. Include screen failure prediction and mitigation strategies.
Process with Claude Code Opus 4.6 (1M token context).]

\subsection{Failure Mode Mitigation}

[PROCESSING INSTRUCTION: Generate 0.75 page with Table 6 covering the 5 failure modes from
sponsor\_04\_operations\_enrollment.md. For each mode, describe the autonomous detection
signal, automated response, and escalation trigger. Include a text-based Gantt timeline
showing the compressed autonomous sponsor timeline from TPP through launch.
Process with Claude Code Opus 4.6 (1M token context).]
